\section{Specify the CRC before writing its loop}
A name such as CRC-32 is insufficient to determine a result. This project uses the
PNG/GZIP polynomial, reflected input processing, initial state \texttt{ffffffff},
and final XOR \texttt{ffffffff}. All hexadecimal state values in this chapter have
exactly 32 significant bits. The right-shifting feedback mask is
\texttt{edb88320}; the corresponding normal polynomial mask is \texttt{04c11db7}.
These masks describe the same polynomial under different bit orientations, but they
cannot be exchanged in the same shift loop.

Keep two values distinct in the interface. The \emph{raw state} is the running
remainder used by an update. The \emph{finalized checksum} is that state XOR
\texttt{ffffffff}. Initialization and finalization belong to the boundary of a
complete message. Neither update routine performs them automatically. For empty
input, the raw state remains \texttt{ffffffff} and its finalized checksum is zero.

A caller starts a message by passing \texttt{ffffffff}, retains each returned raw
state, and complements the last result once. To start an unrelated message, initialize
again. To continue the same message, pass the prior raw result, rather than a printed
checksum or a simulator register left over from a previous run.

\subsection*{Exercise: Identify the boundary operation}
An update returns \texttt{7c231048} after the first byte of a message. State which
value should be passed to the next update and which value would be displayed if that
one byte were the entire message. Explain why displaying a value and continuing a
calculation require different boundary operations.
